% Part: turing-machines
% Chapter: machines-computations
% Section: disciplined-machines

\documentclass[../../../include/open-logic-section]{subfiles}

\begin{document}

\olfileid{tur}{mac}{dis}
\olsection{ماشین‌های منضبط}

\begin{explain}
در بخش \olref[hal]{sec} ماشین‌های تورینگی با یک حالت توقف تعیین‌شدهٔ
$h$ را بررسی کردیم---چنین ماشین‌هایی تضمین می‌کنند که اگر اصلاً
متوقف شوند، در حالت~$h$ متوقف شوند. از این جهت ماشین‌های دارای یک
حالت توقف از آنچه به‌طور کلی برای ماشین‌های تورینگ مجاز می‌دانیم
«منضبط‌ترند». محدودیت‌های دیگری نیز می‌توان بر رفتار ماشین‌های تورینگ
تحمیل کرد. برای نمونه، هنوز ماشین‌های تورینگ را از پاک‌کردن نشانگر
انتهای نوار در خانهٔ~$0$ یا تلاش برای حرکت به چپ از خانهٔ~$0$ منع
نکرده‌ایم. (تعریف ما می‌گوید سر در این حالت صرفاً در خانهٔ~$0$
می‌ماند؛ در تعریف‌های دیگر ماشین متوقف می‌شود.) همچنین گاه مطلوب است
بتوانیم فرض کنیم ماشین تورینگ، اگر اصلاً متوقف شود، در خانهٔ~$1$
متوقف می‌شود.
\end{explain}

\begin{defn}\ollabel{defn:disciplined}
ماشین تورینگ~$M$ \emph{منضبط} است اگر و تنها اگر
\begin{enumerate}
    \item یک حالت توقف یگانه و تعیین‌شدهٔ~$h$ داشته باشد؛
    \item اگر اصلاً متوقف شود، هنگام خواندن خانهٔ~$1$ متوقف شود؛
    \item هرگز نماد $\TMendtape$ را در خانهٔ~$0$ پاک نکند؛ و
    \item هرگز برای حرکت به چپ از خانهٔ~$0$ تلاش نکند.
\end{enumerate}
\end{defn}

\begin{explain}
از پیش گفتیم هر ماشین تورینگ را می‌توان به ماشینی با همان رفتار و یک
حالت توقف تعیین‌شده تبدیل کرد. کافی است حالت تازهٔ~$h$ را بیفزاییم و
برای هر زوج $\tuple{q,\sigma}$ که $\delta$ اصلی در آن تعریف‌نشده است،
دستور $\delta(q,\sigma)=\tuple{h,\sigma,\TMstay}$ را اضافه کنیم. درست است،
گرچه اثباتش ملال‌آور است، که هر ماشین تورینگ~$M$ را می‌توان به ماشین
تورینگ منضبط~$M'$ای تبدیل کرد که روی همان ورودی‌ها متوقف می‌شود و همان
خروجی را تولید می‌کند. برای نمونه، اگر ماشین تورینگ متوقف شود و روی
خانهٔ~$1$ نباشد، می‌توانیم دستورهایی بیفزاییم تا سر به چپ حرکت کند تا
نشانگر انتهای نوار را بیابد، سپس یک خانه به راست برود و متوقف شود.
اندیشیدن دربارهٔ چگونگی رسیدگی به شرط‌های دیگر را به شما واگذار
می‌کنیم.
\end{explain}

\begin{ex}
    در \olref{fig:adder-disc} ماشین جمعِ
    \olref[una]{ex:adder} را به ماشینی منضبط تبدیل می‌کنیم.
    \begin{figure}\[
\begin{tikzpicture}[->,>=stealth',shorten >=1pt,auto,node distance=2.8cm,
                    semithick]
  \tikzstyle{every state}=[fill=none,draw=black,text=black]

  \node[initial,state]         (A)              {$q_0$};
  \node[state]         (B) [right of=A] {$q_1$};
  \node[state]         (C) [below right of=B] {$q_2$};
  \node[state]         (D) [below left of=C] {$q_3$};
  \node[state]         (H) [left of=D] {$h$};

  \path (A) edge node {\TMtrans{\TMblank}{\TMstroke}{\TMstay}} (B)
           edge [loop below] node {\TMtrans{\TMstroke}{\TMstroke}{\TMright}} (A)
        (B) edge [loop below] node {\TMtrans{\TMstroke}{\TMstroke}{\TMright}} (B)
            edge node {\TMtrans{\TMblank}{\TMblank}{\TMleft}} (C)
        (C) edge node {\TMtrans{\TMstroke}{\TMblank}{\TMleft}} (D)
        (D) edge [loop above] node {\TMtrans{\TMstroke}{\TMstroke}{\TMleft}} (D)
        edge node {\TMtrans{\TMendtape}{\TMendtape}{\TMright}} (H);
\end{tikzpicture}
\]\caption{ماشین جمع منضبط}
\ollabel{fig:adder-disc}
\end{figure}
\end{ex}

\begin{prop}\ollabel{prop:disciplined} به ازای هر ماشین تورینگ~$M$،
ماشین تورینگ منضبط~$M'$ای وجود دارد که اگر $M$ با خروجی~$O$ متوقف شود
با خروجی~$O$ متوقف می‌شود، و اگر $M$ متوقف نشود متوقف نمی‌شود.
به‌ویژه، هر تابع $f\colon\Nat^n\to\Nat$ که ماشین تورینگی محاسبه کند
با ماشین تورینگ منضبطی نیز محاسبه‌پذیر است.
\end{prop}

\begin{prob}\label{tur:mac:dis:prob:disc-succ}
ماشین منضبطی به دست دهید که $f(x)=x+1$ را محاسبه کند.
\end{prob}

\begin{prob}\label{tur:mac:dis:prob:copier}
ماشین منضبطی بیابید که با آغاز روی ورودی $\TMstroke^n$، خروجی
$\TMstroke^n\concat\TMblank\concat\TMstroke^n$ را تولید کند.
\end{prob}

\end{document}
